% intro.tex (국문 v13)
활동층은 영구동토 위에서 여름마다 녹고 겨울마다 다시 어는 지표 부근의 지층이다. 활동층 두께(ALT)는 계절 융해의 깊이를 정하며 지반 안정성, 지표 부근 수문, 분해에 노출되는 토양 유기탄소의 양을 좌우한다\cite{Hjort2018,Schuur2015,Biskaborn2019}. ALT는 해마다의 열 조건에 반응하므로, 연도 간 변동은 CALM(Circumpolar Active Layer Monitoring) 같은 장기 관측망이 기록하는 신호의 본질적인 부분이다\cite{Nelson1997,Streletskiy2025}. 그러나 관측은 드물고 공간적으로 모여 있다. 대부분의 측정은 소수의 집중 현장 조사와 장기 격자에서 나오며, 한 번만 방문한 지점이 많다\cite{Moore2025,Veremeeva2025}.

북반구 영구동토 지역에는 약 1 km 해상도의 격자 ALT 제품이 여러 개 있다. ESA CCI Permafrost 같은 과정 기반 추정은 위성 지표 온도로 지중 열전도 방정식을 푼다\cite{Westermann2024,Westermann2025}. 통계 · 기계학습(ML) 제품은 지점 관측을 기후 · 지형 · 토양 예측 변수와 연결한다\cite{Aalto2018,Ran2022a,Liu2024,Wei2026}. 연도별 활용에는 세 가지 제약이 있다. 첫째, 대부분의 제품은 다년 평균 예측 변수나 여러 해의 지점-연도를 합친 자료로 학습되어, 제품 값이 특정 해의 기후와 연결되지 않는다. 둘째, 제품마다 검증 지점과 검증 방식이 달라 보고된 오차를 서로 비교할 수 없다. 예를 들어 Liu 등의 무작위 10겹 RMSE는 21.6 cm\cite{Liu2024}, Wei 등의 지점 하나 제외 RMSE는 60.5 cm\cite{Wei2026}이다. 셋째, 공간적으로 모인 관측에 무작위 교차검증을 적용하면 새 위치에서 ML 모델의 정확도가 과대평가되므로\cite{Roberts2017,Ploton2020,MeyerPebesma2022,Wadoux2021}, 관측이 없는 지역에서의 성능을 판단하기 어렵다.

물리 기반 모델은 다른 접근을 제공한다. Stefan 해는 융해 깊이를 그 계절의 융해 도일(TDD)과 토양의 열적 성질 · 수분에 따라 정해지는 계수 하나로 연결한다\cite{Stefan1891,Riseborough2008,Nelson1997}. 이 해는 계절 융해의 진행을 재현하고 해마다의 기후에 반응하지만, 지역마다 계수 하나를 쓰면 토양 유기물과 지형에 따른 체계적인 잔차가 남는다. 물리 모델이 주된 반응을 맡고 학습기가 물리 모델이 놓친 부분을 보정하는 물리 기반 ML은 여러 환경 분야에서 정확도와 외삽 능력을 높였고\cite{Willard2022,Read2019,Jia2021}, 알래스카에서는 Stefan 모델이 검정 지점에서 랜덤 포레스트보다 정확도를 잘 유지하였다\cite{Gautam2025}. 그러나 기후 입력을 관측 연도와 맞추었을 때 이러한 결합의 성능, 라벨이 없는 지역으로의 전이, 새 지역에 필요한 라벨 수, 같은 셀 · 같은 해에서 공개 제품과의 비교는 밝혀지지 않았다.

이 연구는 지점-연도 설계로 이 문제를 다룬다. 각 ALT 관측은 명시된 해의 관측 지점이며, 각 기후 입력은 그 해의 ERA5-Land 자료로 계산해 관측일까지 누적한다\cite{munozsabater2021_era5_land}. 알래스카, 레나 델타, 캐나다 북서부의 2010--2024년 지점-연도 관측 \nncyAll{}개로 Stefan 모델, 직접 ML, 그리고 재보정한 Stefan 모델의 잔차를 gradient boosting 으로 학습하는 물리 잔차 워크플로를 비교하였다. 제품마다 다른 지점 집합에서 무작위 또는 지점 단위 교차검증을 주로 보고한 기존 ALT 지도 연구\cite{Liu2024,Wei2026,Ran2022a}와 달리, 이 연구는 모든 방법과 제품을 같은 관측과 같은 검증 설계로 채점한다. 이 연구의 기여는 네 가지이다. (i) ALT 라벨과 기후 입력의 시간 불일치를 제거하고 관측 연도 입력의 효과를 정량화하였다. (ii) 물리 잔차 학습을 직접 ML, 물리 입력 ML, 물리 유사라벨과 다섯 가지 검증 설계(지역 전체 제외 포함)에서 비교하였다. (iii) 지역 안 · 지역 간 라벨 수 실험과 외부 CALM 지역 시험으로 표본 효율과 전이성을 정량화하였다. (iv) 모델과 공개 제품 6종을 같은 셀 · 같은 해에서 제품 학습 지점 마스크와 함께 비교하고, 적용 가능 영역과 예측 구간을 갖춘 2011--2024년 연도별 1 km 지도를 제시하였다.
